\documentclass[../main.tex]{subfiles}
\begin{document}
\chapter{Linear Programming}
\section{Formulation}
\begin{definition}[Linear Program]
  A \textit{linear program} has general form:
  \begin{mini*}
    {\vec{x}}{\vec{c}^{\trans}\vec{x}}{}{\label{linearGeneral}}
    \addConstraint{\vec{a}^{\trans}_i \vec{x}}{\geq b_i,\ i \in M_1}
    \addConstraint{\vec{a}^{\trans}_i \vec{x}}{\leq b_i,\ i \in M_2}
    \addConstraint{\vec{a}^{\trans}_i \vec{x}}{= b_i,\ i \in M_3}
    \addConstraint{x_j}{\geq 0,\ j \in N_1}
    \addConstraint{x_j}{\leq 0,\ j \in N_2}
    \addConstraint{x_j}{\text{ free},\ j \in N_3}
  \end{mini*}
\end{definition}
We can in fact write such constraints as a matrix:
\[
  A = \begin{pmatrix}
  \leftarrow \vec{a}^{\trans}_1 \rightarrow \\
  \leftarrow \vec{a}^{\trans}_2 \rightarrow \\
  \cdots \\
  \leftarrow \vec{a}^{\trans}_m \rightarrow
  \end{pmatrix}
\]
\begin{theorem}
  The dual of the problem \cref{linearGeneral} is given by:
\end{theorem}
\begin{proof}
  Adding slack variables, we have:
  \begin{align*}
    L(\vec{x}, \vec{s}, \vec{\lambda}) &= \vec{c}^{\trans} \vec{x} - \sum_{i \in M_1} \lambda_i (\vec{a}^{\trans}_i \vec{x} - s_i - b_i) \\
                     &+ \sum_{i \in M_2} \lambda_i(\vec{a}_i^{\trans} \vec{x} + s_i - b_i) - \sum_{i \in M_3} \lambda_i(\vec{a}^{\trans}_i \vec{x} - b_i) \\
                     &= \sum_{j \in N_1} (c_j - \vec{\lambda}^{\trans} A_j)\vec{x}_j + \sum_{j \in N_2} (c_j - \vec{\lambda}^{\trans} A_j)\vec{x}_j + \sum_{j \in N_3} (c_j - \vec{\lambda}^{\trans} A_j)\vec{x}_j \\
                     &+\sum_{i \in M_1} \lambda_i s_i - \sum_{i \in M_2} \lambda_i s_i + \sum_{i \in M_1 \cup M_2 \cup M_3} \lambda_i b_i
  \end{align*}
  where $A_i$ is the $i$-Th column of $A$.
  Consider the set:
  \[
    \Lambda = \left\{\vec{\lambda}: \begin{cases}
    \lambda_i \geq 0 & \text{ for } i \in M_1 \\
    \lambda_i \leq 0 & \text{ for } i \in M_2 \\
    \lambda_i \text{ free} & \text{ for } i \in M_3 \\
    \vec{\lambda}^{\trans} A_j \leq c_j & \text{ for } j \in N_1 \\
    \vec{\lambda}^{\trans} A_j \geq c_j & \text{ for } j \in N_2 \\
    \vec{\lambda}^{\trans} A_j = c_j & \text{ for } j \in N_3
    \end{cases}\right\}
  \]
  If $\vec{\lambda} \in \Lambda$, then:
  \[
    \min_{\vec{x}, \vec{s} \geq \vec{0}} L(\vec{x}, \vec{s}, \vec{\lambda}) = \vec{b}^{\trans} \vec{\lambda} = g(\vec{\lambda})
  \]
  So we see that the dual is also a linear program.
\end{proof}
\section{Optimality Conditions}
\begin{theorem}
  Let $\vec{x}$ and $\vec{\lambda}$ be feasible solutions to the primal program \cref{linearGeneral}  and its respective $\vec{\lambda}$.
  We then have that these are optimal if and only if:
  \[
    \lambda_i (\vec{a}^{\trans}_i  \vec{x} - b_i) = 0\ \forall i \text{ and } (c_j - \vec{\lambda}^{\trans}A_j)x_j = 0\ \forall j
  \]
\end{theorem}
\begin{proof}
  Define $u_i = \lambda_i(\vec{a}^{\trans}_i \vec{x} - b_i)$ and $v_j = (c_j - \vec{\lambda}^{\trans} A_j)x_j$.
  Since $\vec{x}$ and $\vec{\lambda}$ are primal and dual feasible respectively, we have that $u_i \geq 0\ \forall i$ and $v_j \geq 0\ \forall j$.

  Taking the sum over all $i$ and $j$ respectively:
  \begin{align*}
    \sum_{i} u_i = \vec{\lambda}^{\trans} A \vec{x} - \vec{\lambda}^{\trans} \vec{b} \\
    \sum_{i} v_i = \vec{c}^{\trans} \vec{x} - \vec{\lambda}^{\trans} A\vec{x}
  \end{align*}
  and so:
  \[
    \sum_{i} u_i + \sum_{j} v_j = \vec{c}^{\trans} \vec{x} - \vec{\lambda}^{\trans} \vec{b}
  \]
  By weak duality, we know that $\vec{c}^{\trans} \vec{x} \geq \vec{\lambda}^{\trans} \vec{b}$.
  \begin{proofdirection}{Assume they are optimal}
    So we then have that $\vec{c}^{\trans} \vec{x} = \vec{\lambda}^{\trans} \vec{b}$, we must then have $u_i = 0\ \forall i$ and $v_j = 0\ \forall j$ as this is a sum of non-negative numbers equal to $0$.
  \end{proofdirection}
  \begin{proofdirection}{Assume the equality holds}
    It then follows that $\vec{c}^{\trans} \vec{x} = \vec{\lambda}^{\trans} \vec{b}$ and so $\vec{x}$ and $\vec{\lambda}$ are optimal.
  \end{proofdirection}
\end{proof}
\section{Standard Form for Linear Programs}
Observe that we can always write any linear program as:
\begin{mini*}
  {\vec{x}}{\vec{c}T \vec{x}}{}{}
  \addConstraint{A\vec{x}}{\leq \vec{b}}
\end{mini*}
This is the \textit{general form} for an LP.
We say that an LP is in \textit{standard form} if it is in the form:
\begin{mini*}
  {\vec{x}}{\vec{c}^{\trans} \vec{x}}{}{}
  \addConstraint{A\vec{x}}{= \vec{b}}
  \addConstraint{\vec{x}}{\geq \vec{0}}
\end{mini*}
We then observe that any $x \in \R$ and be expressed as $x = x_{+} - x_{-}$ where $x_{+}, x_{-} \in \R^{+}$.
This can be used to convert any general form problem to a standard form problem.

Starting with:
\begin{mini*}
  {\vec{x}}{\vec{c}T \vec{x}}{}{}
  \addConstraint{A\vec{x}}{\leq \vec{b}}
\end{mini*}
we first add slack variables:
\begin{mini*}
  {\vec{x}}{\vec{c}^{\trans} \vec{x}}{}{}
  \addConstraint{A\vec{x} + \vec{s}}{= \vec{b}}
  \addConstraint{\vec{s}}{\geq 0}
\end{mini*}
and can then use the trick:
\begin{mini*}
  {\vec{x}}{\vec{c}^{\trans}(\vec{x}_{+} - \vec{x}_{-})}{}{}
  \addConstraint{A(\vec{x}_{+} - \vec{x}_{-}) + \vec{s}}{= \vec{b}}
  \addConstraint{\vec{x}_{+}}{\geq 0}
  \addConstraint{\vec{x}_{-}}{\geq 0}
  \addConstraint{\vec{s}}{\geq 0}
\end{mini*}
which is now a standard form problem in $(\vec{x}_{+}, \vec{x}_{-}, \vec{s})$.
\section{Solving Linear Programs}
We first note that minimising $\vec{c}^{\trans} \vec{x}$ is equivalent to maximising $-\vec{c}^{\trans} \vec{x}$.
But $-\vec{c}^{\trans} \vec{x}$ is a convex function, so we are essentially trying to maximise a convex function over some restricted region.
\begin{definition}[Extreme Point]
  Let $C \subseteq \R^{n}$ be a convex set.
  A point $\vec{x} \in C$ is called an \textit{extreme point} if it cannot be written as a convex combination of two distinct points in $C$.

  That is, for all $\vec{y}, \vec{z} \in \C$ and $\delta \in (0, 1)$:
  \[
    \vec{x} = \delta \vec{y} + (1 - \delta) \vec{z} \implies \vec{x} = \vec{y} = \vec{z}
  \]
\end{definition}
\begin{remark}[Intuition]
  A point is extreme if it does not lie on any line segment \textbf{contained within} $C$.

  So if $C$ is a triangle, only its vertices are extreme points, and if $C$ is a circle, then only its boundary points are extreme points.
\end{remark}
\begin{lemma}
  The max of a convex function $f$ on a convex set $C$ occurs at an extreme point of $C$.
\end{lemma}
\begin{proof}
  Suppose $\vec{x}$ is not an extreme point.
  Thus there is $\vec{y}, \vec{z} \in C$ and $\delta \in (0, 1)$ such that:
  \[
    f(\vec{x}) \leq \delta f(\vec{y}) + (1 - \delta)f(\vec{z}) \leq \max\{f(\vec{y}), f(\vec{z})\}
  \]
\end{proof}
In a linear program, the constraint set is always a polytope.
Such regions always have finitely many extreme points and so if we can find a list of these extreme points it becomes easier to maximise $-\vec{c}^{\trans} \vec{x}$.
\subsection{Basic Solutions and Basic Feasible Solutions}
Recall that we wish to solve:
\begin{mini*}
  {\vec{x}}{\vec{c}^{\trans} \vec{x}}{}{}
  \addConstraint{A\vec{x}}{= \vec{b}}
  \addConstraint{\vec{x}}{\geq \vec{0}}
\end{mini*}
where:
\[
  A = \begin{pmatrix}
  \leftarrow \vec{a}^{\trans}_1 \rightarrow \\
  \leftarrow \vec{a}^{\trans}_2 \rightarrow \\
  \cdots \\
  \leftarrow \vec{a}^{\trans}_m \rightarrow
  \end{pmatrix}
\]
We will make three assumptions:
\begin{enumerate}
  \item The rows of $A$ are linearly independent.
    This can in fact be made without loss in generality as if two rows are not linearly independent then we have a redundant constraint so we can simply remove one of the rows.
  \item Every set of $m$ columns of $A$ is linearly independent.
    This cannot be made without loss in generality, but for most applications, it will typically hold.
\end{enumerate}
For our final assumption, we will need some additional definitions.
\begin{definition}[Basic Solution]
  We call a point $\vec{x}$ that satisfies $A\vec{x} = \vec{b}$ and has at most $m$ non-zero entries a \textit{basic solution}.
\end{definition}
\begin{definition}[Basic Feasible Solution]
  We call a basic solution $\vec{x}$ that satisfies $\vec{x} \geq \vec{0}$ a \textit{basic feasible solution}.
\end{definition}
Suppose we choose indices $\{B(1), \ldots, B(m)\}$ such that $x_i = 0$ for all $i \notin \{B(1), \ldots, B(m)\}$, i.e. we have a set of indices we know have $x_i = 0$ and a set of indices we know have $x_i \neq 0$:
\[
  A \vec{x} = \sum_{i = 1}^{n} \vec{A}_i x_i = \sum_{i = 1}^{m} \vec{A}_{B(i)}\vec{x}_{B(i)} = B \vec{x}_{B}
\]
where $\vec{A}_i$ is the $i$-th column of $A$ and:
\[
  B = \begin{pmatrix}
  \uparrow & & \uparrow  \\
  \vec{A}_{B(1)} & \cdots & \vec{A}_{B(m)}  \\
  \downarrow &  & \downarrow \\
  \end{pmatrix}
 \text{ and } \vec{x}_{B} = \begin{pmatrix}
  x_{B(1)} \\
  \vdots \\
  x_{B(m)} \\
  \end{pmatrix}
\]
Ideally for this problem, we want $A\vec{x} = \vec{b} \iff B\vec{x}_{B} = \vec{b} \iff \vec{x}_{B} = B^{-1} \vec{b}$.
\begin{definition}[Basic Variables, Matrix, and Columns]
  We call:
  \begin{itemize}
    \item $\vec{x}_{B(1)}, \ldots, \vec{x}_{B(m)}$ the \textit{basic variables},
    \item $B$ the \textit{basis matrix},
    \item $\vec{A}_{B(1)}, \ldots, \vec{A}_{B(m)}$ the \textit{basic columns}.
  \end{itemize}
\end{definition}
For every choice of $B$ (of which there are $\binom{n}{m}$ such choices for the columns), we obtain a basic solution.
So if $B^{-1}\vec{b} \geq 0$, then that the basis matrix $B$ leads to a basic feasible solution.

We can now define our third assumption:
\begin{enumerate}
  \setcounter{enumi}{2}
  \item Every basic feasible solutions is \textit{non-degenerate}.
    That is, it has \textbf{exactly} $m$ non-zero entries:
    \[
      B^{-1} \vec{b} \geq 0 \implies B^{-1} \vec{b} > 0
    \]
    This also cannot be made without loss in generality but is required to ensure that algorithms later in this chapter do not get stuck.
\end{enumerate}
\begin{theorem}
  A vector $\vec{x}$ is a basic feasible solution if and only if it is an extreme point of the set:
  \[
    \mathcal{X}(\vec{b}) = \{\vec{x}: A \vec{x} = \vec{b}, \vec{x} \geq 0\}
  \]
\end{theorem}
\begin{remark}[Rudimentary Algorithm for Solving a L.P.]
  This theorem means that we can solve an L.P. by just checking each of the basic feasible solutions, that is:
  \begin{itemize}
    \item List all $\binom{n}{m}$ basic solutions.
    \item Filter these basic solutions to just the basic feasible solutions.
    \item Evaluate $\vec{c}^{\trans} \vec{x}$ for each of these basic feasible solutions and choose the best one.
  \end{itemize}
\end{remark}
\begin{proof}
  \begin{proofdirection}{Assume $\vec{x}$ is a basic feasible solution}
    Suppose that $\vec{x}$ can be written as:
    \[
      \vec{x} = (1 - \delta)\vec{y} + \delta \vec{z}
    \]
    for some $\vec{y}, \vec{z} \in \mathcal{X}(\vec{b})$ and $\delta \in (0, 1)$.

    If $x_i = 0$, then:
    \[
      x_i = 0 = (1 - \delta)y_i + \delta z_i
    \]
    but we also have that $y_i, z_i \geq 0$ as $\vec{y}, \vec{z} \in \mathcal{X}(\vec{b})$ and so $y_i = z_i = 0$.
    Since $\vec{x}$ is a basic feasible solution, $\vec{y}$ and $\vec{z}$ must be basic solutions corresponding to the same basis as $\vec{x}$.
    Thus $\vec{x} = \vec{y} = \vec{z} = B^{-1} \vec{b}$ and so $\vec{x}$ is an extreme point of $\mathcal{X}(\vec{b})$.
  \end{proofdirection}
  \begin{proofdirection}{Assume $\vec{x}$ is not a basic feasible solution}
    We wish to show that $\vec{x}$ is not an extreme point of $\mathcal{X}(\vec{b})$.
    If $\vec{x} \notin \mathcal{X}(\vec{b})$, then obviously it is not an extreme point  of $\mathcal{X}(\vec{b})$ and so we assume that $\vec{x} \in \mathcal{X}(\vec{b})$.

    Since $\vec{x}$ is not a basic feasible solution, it must have $r > m$ positive entries.
    Suppose that $x_{i_1} > 0, x_{i_2} > 0, \ldots, x_{i_r} > 0$.
    Because we have chosen more columns than the rank $m$, the set of columns $\{\vec{A}_{i_1}, \ldots, \vec{A}_{i_r}\}$ must be linearly dependent.
    By definition this means that $\exists w_{i_1}, \ldots, w_{i_r}$ not all 0 such that:
    \[
      w_{i_1} \vec{A}_{i_1} + \cdots + w_{i_r} \vec{A}_{i_r} = \vec{0}
    \]
    Construct a vector $\vec{w}$ such that $\vec{w}$ is $0$ everywhere except that $i_1, \ldots, i_r$ where $w(i_1) = w_{i_1}, \ldots, w(i_r) = w_{i_r}$.

    Now observe that $A\vec{w} = \vec{0}$ because:
    \[
      A \vec{w} = \vec{A}_{i_1} w_{i_1} + \cdots + \vec{A}_{i_r} w_{i_r}
    \]
    Now consider two new points:
    \begin{align*}
      \vec{y} = \vec{x} + \varepsilon \vec{w} \\
      \vec{z} = \vec{x} - \varepsilon \vec{w}
    \end{align*}
    for some small $\varepsilon > 0$.

    For small enough $\varepsilon$, $\vec{y}, \vec{z} \in \mathcal{X}(\vec{b})$.
    Thus we can write $\vec{x}$ as:
    \[
      \vec{x} = \frac{1}{2}\vec{y} + \frac{1}{2}\vec{z}
    \]
    for $\vec{y}, \vec{z} \in \mathcal{X}(\vec{b})$ and so $\vec{x}$ cannot be an extreme point of $\mathcal{X}(\vec{b})$.
  \end{proofdirection}
\end{proof}
\subsection{Towards the Simplex Method}
When the L.P. is in standard form, the optimality conditions are:
\begin{enumerate}
  \item \textbf{Primal feasibility --} $A\vec{x} = \vec{b}$, $\vec{x} \geq 0$.
  \item \textbf{Dual feasibility --} $A^{\trans} \vec{\lambda} \leq \vec{c}$
  \item \textbf{Complementary Slackness --} $\vec{x}^{\trans}(\vec{c} - A^{\trans} \vec{\lambda}) = 0$
\end{enumerate}
For any basic feasible solution, the complementary slackness equation reduces to:
\[
  \vec{x}_{B}^{\trans}(\vec{c}_{B} - B^{\trans} \vec{\lambda}) = 0
\]
Since $\vec{x}_{B} > \vec{0}$ by the non-degeneracy assumption, we must have:
\[
  \vec{c}_{B} = B^{T} \vec{\lambda} \implies \vec{\lambda} = (B^{\trans})^{-1} \vec{c}_{B}
\]
If $A^{\trans}(B^{\trans})^{-1}\vec{c}_{B} \leq \vec{c}$, then $\vec{\lambda}$ is dual feasible and so $\vec{x}$ is optimal.

So if we check $A^{\trans}(B^{\trans})^{-1}\vec{c}_{B} \leq \vec{c}$ at every basic feasible solution, any solution satisfying this must be optimal.

\begin{definition}[Vector of reduced costs]
  We define the \textit{vector of reduced costs} $\overline{\vec{c}}$ as:
  \[
    \overline{\vec{c}}^{\trans} = \vec{c}^{\trans} - \vec{c}^{\trans}_{B} B^{-1} A
  \]
\end{definition}
So for a BFS $\vec{x}$, if $\overline{\vec{c}} \geq \vec{0}$, then $\vec{x}$ must be optimal.
\section{Simplex Algorithm}
\subsection{Theory}
Suppose we are at the point:
\[
  \vec{x} = (x_{B(1)}, \ldots, x_{B(m)}, 0, \ldots, 0)^{\trans}
\]
Suppose for some $j^{*}$, $\overline{c}_{j^{*}} < 0$, i.e., $\vec{x}$ is not optimal.
Suppose now that we choose some $j \notin \{B(1), \ldots, B(m)\}$ and we decide to make $x_j > 0$ whilst keeping all non basic $x_i = 0$:
\[
  \vec{x}' = (x_{B(1)}, \ldots, x_{B(m)}, 0, \ldots, 0)^{\trans} + t (d_{B(1)}, \ldots, d_{B(m)}, 0, \ldots, 1, \ldots, 0)^{\trans}
\]
where the $1$ is in the $j^{*}$-th position.

We need $A\vec{x}' = \vec{b}$ and so:
\begin{align*}
  A(\vec{x} + t(d_{B(1)}, \ldots, d_{B(m)}, 0, \ldots, 1, \ldots, 0)^{\trans}) &= \vec{b} \\
  \implies \cancel{\vec{b}} + tA(d_{B(1)}, \ldots, d_{B(m)}, 0, \ldots, 1, \ldots, 0)^{\trans} &= \cancel{\vec{b}} \\
  \implies tA(d_{B(1)}, \ldots, d_{B(m)}, 0, \ldots, 1, \ldots, 0)^{\trans} &= \vec{0}
\end{align*}
If this $\vec{x}'$ is still feasible then:
\begin{align*}
  A(d_{B(1)}, \ldots, d_{B(m)}, 0, \ldots, 1, \ldots, 0)^{\trans} = \vec{0} &\iff B \vec{d}_{B} + \vec{A}_j = 0\ \forall j \\
                                                                            &\implies \vec{d}_B = - B^{-1} \vec{A}_j\ \forall j
\end{align*}
We now check the cost at $\vec{x}'$:
\begin{align*}
  \vec{c}^{\trans} \vec{x}' &= \vec{c}^{\trans}(\vec{x} + t(\vec{d}_B, 0, \ldots, 1, \ldots, 0)^{\trans}) \\
                            &= \vec{c}^{\trans} \vec{x} + t(\vec{c}^{\trans}_{B} \vec{d}_B + c_j) \\
                            &= \vec{c}^{\trans} \vec{x} + t(c_j - \vec{c}^{\trans}_{B} B^{-1}\vec{A}_j) \\
                            &= \vec{c}^{\trans} \vec{x} + tc_j
\end{align*}
We then calculate:
\[
  - \frac{x_{B(1)}}{d_{B(1)}}, \ldots, - \frac{x_{B(m)}}{d_{B(m)}}
\]
and find the smallest positive one.

Let us consider how is $x_{B(1)}$ changes with $t$:
\begin{align*}
  x_{B(1)} + t d_{B(1)} \geq 0 \iff x_{B(1)} \geq - t d_{B(1)}
\end{align*}
But if $d_{B(1)} < 0$, then the $x_{B(1)}$ could become negative and so we would need to ensure that $t \leq -\frac{x_{B(1)}}{d_{B(1)}}$.

If $d_{B(i)} \geq 0$, then we do not need to worry about leaving the constrained set in this direction, only those $B(i)$ with $d_{B(i)} < 0$ can cause constraints issues, hence why we take:
\[
  t^{*} = \min\left\{- \frac{x_{B(1)}}{d_{B(1)}}, \ldots, - \frac{x_{B(m)}}{d_{B(m)}}\right\}
\]
What can we say about $\vec{x}'$ when we take $t = t^{*}$?
\[
  \vec{x}' = \vec{x} + t^{*} \vec{d} \text{ where } \vec{d} = (\vec{d}_B, 0, \ldots, 1, \ldots, 0)^{\trans}
\]
\begin{proposition}
  $\vec{x}' = \vec{x} + t^{*} \vec{d}$ as defined above is a basic feasible solution.
\end{proposition}
\begin{proof}
  We already have that $A \vec{y} = \vec{0}$ and $\vec{y} \geq \vec{0}$ as our choice of $t^{*}$ ensures that $\vec{y}$ is feasible.

  Because of our choice of $t^{*}$, \textbf{exactly 1} amongst $x_{B(1)}', \ldots, x_{B(m)}'$ has become 0 and $x_{j^{*}}' > 0$.
  We know that there is exactly 1 which is $0$ because of our assumption that all of our basic feasible solutions are non-degenerate.

  Thus $y$ has exactly $m$ non-zero entries and so is a basic feasible solution.
\end{proof}
So we have figured out a method for starting at a basic feasible solution to find a new basic feasible solution where the reduced cost is strictly less than our original reduced cost.
\subsection{Algorithm}
We start at a basic feasible solution $\vec{x}$ and construct a \textit{simplex tableau}:
\begin{center}
\begin{tabular}{c|c|c|c|c}
$\vec{c}^{\trans}_{B} \vec{x}_{B}$ & $\overline{c}_1$ & $\overline{c}_2$ & $\cdots$ & $\overline{c}_n$ \\
\hline
$x_{B(1)}$ & $\uparrow$ & $\uparrow$ &  & $\uparrow$ \\
$\vdots$ & $B^{-1}\vec{A}_1$ & $B^{-1} \vec{A}_2$ & $\cdots$ & $B^{-1} \vec{A}_n$ \\
$x_{B(m)}$ & $\downarrow$ & $\downarrow$ &  & $\downarrow$
\end{tabular}
\end{center}
%TODO Review
\begin{enumerate}
  \item Convert the general form problem into a standard form problem by introducing slack variables:
    \begin{enumerate}
      \item Add a positive slack variable for $\leq$ and add it as a component of $\vec{x}$.
      \item Add a negative slack variable for $\geq$ and add it as a component of $\vec{x}$.
    \end{enumerate}
  \item Convert the problem into the vectors $\vec{b}$, $\vec{c}$ and matrices $A$, $B$:
    \begin{enumerate}
      \item $A$ is the matrix of coefficients of the variables in the constraints.
      \item $\vec{b}$ is the vector in the constraint $A\vec{x} = \vec{b}$.
      \item $\vec{c}$ is the vector of coefficients of the variables in the objective function.
      \item
    \end{enumerate}
  \item Start the simplex method with a basic feasible solution by setting all initial variables to 0 and all slack variables equal to the $b_i$ in the corresponding constraint.
    This step only works for $\vec{b} \geq \vec{0}$, we can in fact deal with the case where we do not have $\vec{b} \geq 0$ but we will not cover this in this course.
    The matrix $B$ is then the matrix from the last columns corresponding to only the slack variables, in the first step, this will always be an identity matrix.
  \item Fill in the simplex tableau.
  \item Identify the direction to move in:
    \begin{enumerate}
      \item If all entries in the top row are non-negative, i.e. $\overline{\vec{c}} \geq \vec{0}$, then we are done.
      \item If not, pick the column corresponding to a negative entry in the top row, if there is more than one, pick any of them, call this column $j$.
      \item For each value in the $j$-th column, divide the element in the first column and same row.
        Pick the row corresponding to the minimal ratio, call this row $i$.
    \end{enumerate}
  \item Perform row operations so that the $j$-th column has 1 in the $i$-th row and 0 in all others.
  \item Return to step \textbf{(iv)} with a new simplex tableau until the algorithm terminates.
\end{enumerate}
\end{document}
